\chapter{Metodologia di sviluppo e tecnologie adottate}

\section{Metodologia di sviluppo}

Lo sviluppo del sistema \textit{JarFin} è stato condotto adottando una metodologia Agile, con l’obiettivo di favorire un processo di progettazione incrementale, flessibile e orientato al cambiamento. In particolare, il progetto ha seguito i principi di \textbf{SCRUM}, integrati con l’approccio \textbf{Agile Model Driven Development (AMDD)}, al fine di bilanciare la necessità di una progettazione architetturale solida con i vantaggi dell’iterazione continua e del feedback costante.

SCRUM è stato utilizzato come framework di riferimento per l’organizzazione del lavoro, la pianificazione delle attività e la gestione delle iterazioni di sviluppo. Il ciclo di vita del progetto è stato articolato in sprint di durata indicativa pari a due settimane, ciascuno caratterizzato da obiettivi ben definiti e dalla produzione di un incremento funzionale del sistema. Questa impostazione ha consentito di monitorare costantemente l’avanzamento del progetto e di adattare le scelte progettuali in base alle esigenze emerse durante lo sviluppo \cite{scrum}.

Parallelamente, l’approccio AMDD ha guidato le attività di modellazione e progettazione. Invece di definire in modo esaustivo l’intera architettura nelle fasi iniziali, la modellazione è stata condotta in maniera incrementale, producendo modelli essenziali e sufficienti a supportare le decisioni architetturali di ciascuna iterazione. In particolare, sono stati realizzati modelli di dominio e diagrammi UML (use case, componenti, classi e sequenza) per i casi d’uso più rilevanti, successivamente raffinati nel corso dello sviluppo \cite{amdd}. Questo approccio ha permesso di evitare un’eccessiva rigidità progettuale, mantenendo al contempo una visione coerente dell’architettura complessiva del sistema.

\section{Organizzazione delle iterazioni}

Il ciclo di sviluppo di JarFin è stato articolato in una serie di iterazioni successive, ciascuna focalizzata su specifici obiettivi funzionali e architetturali:

\begin{itemize}
\item \textbf{Iterazione 0 – Envisioning e Analisi}: definizione del dominio applicativo, raccolta e analisi dei requisiti funzionali e non funzionali, identificazione dei principali casi d’uso e progettazione dell’architettura di alto livello;
\item \textbf{Iterazione 1 – Core Services}: progettazione e implementazione dei servizi fondamentali del sistema, con particolare riferimento al servizio di accounting e alla gestione della persistenza dei dati;
\item \textbf{Iterazione 2 – Analytics e Algoritmi}: sviluppo del servizio di analytics e implementazione degli algoritmi per l’elaborazione e l’analisi dei dati finanziari;
\item \textbf{Iterazione 3 – Integrazione e Interazione}: integrazione del modulo di Natural Language Understanding, dell’API Gateway e dell’interfaccia web.
\end{itemize}

Questa suddivisione ha consentito di affrontare progressivamente la complessità del sistema, validando le scelte progettuali a ogni iterazione e riducendo il rischio di errori architetturali nelle fasi avanzate dello sviluppo.

\section{Architettura generale del sistema}

L’architettura di JarFin è basata su un paradigma a microservizi, in cui il sistema è scomposto in un insieme di componenti indipendenti, ciascuno responsabile di una specifica area funzionale. Tale scelta architetturale è motivata dall’esigenza di garantire modularità, scalabilità e manutenibilità, oltre a facilitare l’evoluzione futura del sistema.

Ogni microservizio è progettato come un’unità autonoma, dotata di una propria logica applicativa e, ove necessario, di un proprio livello di persistenza. La comunicazione tra i servizi avviene tramite API REST e scambio di messaggi in formato JSON, secondo i principi dello stile architetturale REST \cite{fieldingRest}.

L’adozione di un’architettura distribuita consente inoltre di isolare le responsabilità dei singoli componenti, riducendo l’impatto delle modifiche e migliorando la resilienza complessiva del sistema in caso di malfunzionamenti localizzati.

\section{Tecnologie adottate}

La scelta dello stack tecnologico è stata guidata dalla necessità di utilizzare strumenti robusti, ampiamente supportati dalla community e adatti allo sviluppo di sistemi distribuiti.

\subsection{Spring Boot}

Per l’implementazione dei microservizi è stato utilizzato il framework \textbf{Spring Boot}, che fornisce un ambiente di sviluppo semplificato per la realizzazione di applicazioni Java basate su architettura a servizi. Spring Boot consente di ridurre significativamente la configurazione manuale, favorendo lo sviluppo rapido di applicazioni standalone e facilmente deployabili \cite{springboot}.

Il framework supporta nativamente la creazione di API REST, l’integrazione con sistemi di persistenza e la gestione della configurazione dei servizi, risultando particolarmente adatto allo sviluppo di architetture a microservizi.

\subsection{Database PostgreSQL}

La persistenza dei dati è affidata a un database relazionale \textbf{PostgreSQL}, scelto per la sua affidabilità, conformità agli standard SQL e supporto a funzionalità avanzate di gestione dei dati. PostgreSQL garantisce l’integrità dei dati finanziari e supporta operazioni di interrogazione efficienti, fondamentali per le funzionalità di analisi offerte dal sistema.

\subsection{Modulo di Natural Language Understanding}

Per l’elaborazione del linguaggio naturale è stato sviluppato un modulo di Natural Language Understanding implementato in Java come microservizio dedicato. Il sistema utilizza tecniche di analisi lessicale e \textit{pattern matching} per identificare intenti e parametri rilevanti a partire dalle frasi dell’utente, consentendo l’esecuzione di operazioni finanziarie tramite comandi in linguaggio naturale.

La scelta di una soluzione interna, anziché basata su servizi cloud esterni, consente di ridurre la latenza, mantenere il controllo sui dati sensibili dell’utente e garantire una maggiore integrazione con l’architettura complessiva del sistema.

\subsection{Containerizzazione e ambiente di esecuzione}

Per facilitare il deployment e garantire la riproducibilità dell’ambiente di esecuzione, il sistema è stato containerizzato mediante \textbf{Docker}. L’utilizzo di container consente di isolare i singoli microservizi e le relative dipendenze, semplificando la gestione dell’infrastruttura e favorendo la portabilità del sistema tra ambienti di sviluppo e produzione.

\section{Toolchain di progetto}

Lo sviluppo del sistema \textit{JarFin} è stato supportato da una toolchain articolata, composta da strumenti dedicati alla modellazione, allo sviluppo, al versioning, al testing e all’analisi della qualità del codice.  
La scelta degli strumenti è stata guidata dall’obiettivo di garantire un processo di sviluppo strutturato, tracciabile e conforme alle buone pratiche dell’ingegneria del software.

\subsection{Strumenti di modellazione}

Per la modellazione del sistema è stato utilizzato \textbf{UMLet}, uno strumento leggero per la creazione di diagrammi UML. UMLet è stato impiegato per la realizzazione dei principali artefatti di modellazione previsti dall’approccio AMDD, in particolare:
\begin{itemize}
\item diagrammi dei casi d’uso, per rappresentare le funzionalità offerte dal sistema e le interazioni con gli attori;
\item diagrammi delle classi, per descrivere la struttura statica dei principali componenti software;
\item diagrammi dei componenti, per modellare l’architettura a microservizi;
\item diagrammi di sequenza, per analizzare il comportamento dinamico del sistema nei casi d’uso critici.
\end{itemize}

L’utilizzo di UMLet ha consentito di produrre modelli essenziali e facilmente modificabili, in linea con i principi della modellazione incrementale promossi dall’Agile Model Driven Development.

\subsection{Versioning e Project Management}

Il controllo di versione del codice sorgente e la gestione del progetto sono stati affidati a \textbf{GitHub}.  
Git è stato utilizzato come sistema di versioning distribuito, permettendo di tracciare l’evoluzione del codice, gestire rami di sviluppo separati e mantenere uno storico completo delle modifiche.

GitHub è stato inoltre impiegato come strumento di supporto al project management, sfruttando:
\begin{itemize}
\item repository centralizzato per il codice sorgente;
\item Kanban board per la pianificazione delle attività e il tracciamento delle iterazioni;
\item Pull Request per la revisione del codice e l’integrazione controllata delle funzionalità.
\end{itemize}

Questa impostazione ha favorito un processo di sviluppo ordinato e trasparente, coerente con la metodologia SCRUM adottata.

\subsection{Sviluppo backend}

Il backend del sistema è stato sviluppato utilizzando il linguaggio \textbf{Java} e il framework \textbf{Spring Boot}.  
Spring Boot è stato scelto per la sua capacità di semplificare lo sviluppo di applicazioni basate su architettura a microservizi, riducendo la configurazione manuale e favorendo la realizzazione di servizi REST standalone.

Ogni microservizio di JarFin è implementato come applicazione Spring Boot indipendente, responsabile di una specifica area funzionale del sistema, in linea con i principi di separazione delle responsabilità e modularità.

\subsection{Testing delle API REST}

Per la verifica manuale e il testing degli endpoint REST è stato utilizzato \textbf{Postman}.  
Lo strumento ha permesso di:
\begin{itemize}
\item testare le operazioni CRUD esposte dai microservizi;
\item validare il corretto formato delle richieste e delle risposte JSON;
\item simulare scenari di utilizzo reali durante le fasi di sviluppo e integrazione.
\end{itemize}

L’uso di Postman ha facilitato l’individuazione precoce di errori di comunicazione tra i servizi e ha supportato la fase di debugging dell’API Gateway.

\subsection{Unit Testing}

Il testing automatico della logica applicativa è stato realizzato tramite \textbf{JUnit 5}, framework di riferimento per il testing in ambiente Java.  
JUnit è stato utilizzato per definire test unitari mirati a verificare il corretto funzionamento dei metodi di business logic, in particolare nei servizi di accounting, analytics e NLU.

L’adozione del testing unitario ha contribuito a migliorare l’affidabilità del sistema e a ridurre il rischio di regressioni durante l’evoluzione del codice.

\subsection{Analisi della copertura del codice}

Per l’analisi della copertura dei test è stato utilizzato il plugin \textbf{EclEmma}, integrato nell’ambiente di sviluppo.  
EclEmma consente di misurare dinamicamente la percentuale di codice eseguito durante l’esecuzione dei test, fornendo indicatori utili per valutare l’efficacia della suite di test.

Le metriche di copertura hanno supportato l’identificazione delle porzioni di codice meno testate, guidando il miglioramento progressivo dei test unitari.

\subsection{Analisi statica del codice}

La qualità strutturale del codice è stata analizzata mediante \textbf{STAN4J}, strumento di analisi statica per applicazioni Java.  
STAN4J è stato utilizzato per:
\begin{itemize}
\item valutare metriche di complessità e dipendenze tra classi;
\item individuare potenziali criticità architetturali;
\item supportare la verifica del rispetto dei principi di progettazione software.
\end{itemize}

L’analisi statica ha contribuito a mantenere un’architettura pulita e coerente con i principi di modularità e manutenibilità, aspetti fondamentali in un sistema a microservizi.

\subsection{Considerazioni sulla toolchain}

L’integrazione degli strumenti descritti ha permesso di coprire l’intero ciclo di vita del software, dalla modellazione iniziale alla verifica della qualità del codice.  
La toolchain adottata rappresenta una soluzione completa e coerente con gli obiettivi del progetto, fornendo un supporto efficace sia allo sviluppo incrementale sia alla validazione tecnica del sistema.